\chapter{Series de Fourier (Tiempo discreto)}

\section{Definiciones}

\subsection{Señales periódicas}

Recordemos que una funci\'{o}n $x[n]$ es peri\'{o}dica si y solo si existe un n\'{u}mero entero $N > 0$
que cumple la condici\'{o}n
\begin{equation}
x\left[ n\right] =x\left[ n+N\right],
\end{equation}
para todo $n$. 

Por ejemplo, 
\begin{eqnarray}
x[n] &=&\cos (\frac{2 \, \pi }{N} \, n) = \cos(\frac{2 \, \pi}{N} \, (n + N)) = \cos (\frac{2 \, \pi}{N} \, n + 2 \, \pi),\\
x\left[ n\right] &=&\sin (\frac{2 \, \pi }{N} \, n) = 
\sin(\frac{2 \pi}{N} (n+N)) = \sin (\frac{2 \pi}{N} \, n + 2 \, \pi), \\
x\left[ n\right] &=&
e^{j\frac{2 \, \pi }{N} \, n} = 
e^{j \, \frac{2 \, \pi}{N} \, (n + N)} = 
e^{j \, \frac{2 \, \pi}{N} + j \, 2 \, \pi} = 
e^{j \, \frac{2 \, \pi}{N}} \, e^{j \, 2 \, \pi}. 
\end{eqnarray}

Recordemos que $e^{j \, 2 \, \pi} = cos(2 \, \pi) + j \, sin(2 \, \pi) = 1$.

\subsection{Se\~{n}ales peri\'{o}dicas arm\'{o}nicamente relacionadas}

Las se\~ {n}ales per\'{o}dicas arm\'{o}nicamente relacionadas cumplen con la condici\'{o}n
\begin{equation}
\phi _{k}\left[ n\right] =e^{j \, k \, \omega \, n}=e^{j \, k \, \frac{2 \, \pi }{N} \, n}, \qquad \forall k=0,\pm 1,\pm 2,...  \label{NHar}
\end{equation}
o la condici\'{o}n
\begin{equation}
\alpha _{k}\left[ n\right] =\cos \left( k \, \omega \, n\right) =
\cos \left(k \, \frac{2 \, \pi }{N} \, n\right), \qquad \forall k=0, 1, 2,..., N - 1, 
\end{equation}
o la condici\'{o}n
\begin{equation}
\beta _{k}\left[ n\right] =\sin \left( k \, \omega \, n\right) =
\sin \left( k \, \frac{2 \, \pi }{N} \, n\right), \qquad \forall k=0, 1, 2,..., N - 1.
\end{equation}
En otras palabras, son funciones peri\'{o}dicas cuyas frecuencias son m\'{u}ltiplos enteros de una frecuencia fundamental.

\begin{enumerate}
\item  Cada una de estas se\~{n}ales tiene una frecuencia fundamental $\omega $ y 
per\'{\i}odo fundamental $N$.

\item  La arm\'{o}nica $k-$ésima, para $k \neq 0$, es peri\'{o}dica con \textbf{per\'{\i}odo }$\frac{N}{k}$.

\item  Solamente existen $N$ se\~{n}ales distintas en el conjunto descripto
por (\ref{NHar}) o sea 
\begin{equation}
\phi _{k}\left[ n\right] =\phi _{k+rN}\left[ n\right].
\end{equation}
Para el caso de las exponenciales complejas podemos demostrar
\begin{equation}
e^{j \, (k + r \, N) \, \frac{2 \, \pi}{N}} =
e^{j \, k \, \frac{2 \, \pi}{N} + r \, N \, \frac{2 \, \pi}{N}} =
e^{j \, k \, \frac{2 \, \pi}{N} + r \, 2 \, \pi} =
e^{j \, k \, \frac{2 \, \pi}{N}} \, e^{r \, 2 \, \pi} = 
e^{j \, k \, \frac{2 \, \pi}{N}}.
\end{equation}

\end{enumerate}

Por ejemplo, para $N = 4$ tenemos las armónicas
\begin{equation}
 x_{0}[n] = 1, 
 \qquad x_{1}[n] = e^{j \, 1 \, \frac{2 \, \pi}{4} \, n}, 
 \qquad x_{2}[n] = e^{j \, 2 \, \frac{2 \, \pi}{4} \, n}, 
 \qquad x_{3}[n] = e^{j \, 3 \, \frac{2 \, \pi}{4} \, n}.
\end{equation}
Observemos las igualdades
\begin{equation}
 x_{k}[n] = e^{j \, k \, \frac{2 \, \pi}{4} \, n} = x_{k}[n + 4] = e^{j \, k \, \frac{2 \, \pi}{4} \, (n + 4)} =
 e^{j \, k \, \frac{2 \, \pi}{4} \, n} \, e^{j \, k \, \frac{2 \, \pi}{4} \, 4} =
  e^{j \, k \, \frac{2 \, \pi}{4} \, n},
\end{equation}
y
\begin{equation}
 x_{4}[n] = e^{j \, 4 \, \frac{2 \, \pi}{4} \, n} = e^{j \, 2 \, \pi \, n} = 1 = x_{0}[n].
\end{equation}


\section{Serie de Fourier de tiempo discreto}

\subsection{Síntesis}

La serie de Fourier de tiempo discreto es una sumatoria de $N$ arm\'{o}nicas complejas de per\'{i}odo $N$, y la podemos escribir as\'{\i}
\begin{equation}
x\left[ n\right] =
\sum_{k=0}^{N-1}a_{k} \, \phi _{k}\left[ n\right]
=
\sum_{k=0}^{N-1}a_{k} \, e^{j \, k \, \omega \, n}
=
\sum_{k=0}^{N-1}a_{k} \, e^{j \, k \, \left( \frac{2 \, \pi }{N}\right) \, n}.
\end{equation}
\begin{enumerate}
\item La sumatoria debe incluir solamente $N$ arm\'{o}nicas. 
\item El valor inicial de $k$ puede ser cualquiera.
\end{enumerate}

\subsection{Análisis, o el cálculo de coeficientes de la serie de Fourier}

\begin{enumerate}
\item Comenzamos escribiendo la serie de Fourier de una se\~{n}al peri\'{o}dica de tiempo discreto 
\begin{equation}
        x\left[ n\right] =
                \sum_{k=0}^{N-1}a_{k} \, e^{j \, k \, \frac{2 \, \pi }{N} \, n}.
\end{equation}
Multiplicando la serie por $e^{-j \, r \, \frac{2 \, \pi }{N} \, n}$ obtenemos 
\begin{equation}
        x\left[ n\right] \, e^{-j \, r \, \frac{2 \, \pi }{N} \, n}=
                \sum_{k=0}^{N-1}a_{k} \, e^{j \, k \, \frac{2 \, \pi }{N} \, n} \, e^{-j \, r \, \frac{2 \, \pi }{N} \, n},
\end{equation}

\item Sumamos ambos lados de la igualdad desde $n=0$ hasta $N-1$
\begin{equation}
        \sum_{n=0}^{N-1}x\left[ n\right] \, e^{-j \, r \, \frac{2 \, \pi }{N} \, n}
        =
        \sum_{n=0}^{N-1}\sum_{k=0}^{N-1}a_{k} \, e^{j \, \left( k-r\right) \, \frac{2 \, \pi }{N} \, n},
\end{equation}

\item Invertimos la posici\'{o}n de los signos de sumatoria a la derecha del signo $=$ y dado que $a_{k}$ no depende de $n$ se obtiene 
\begin{equation}
        \sum_{n=0}^{N-1}x\left[ n\right] \, e^{-j \, r \, \frac{2 \, \pi }{N} \, n}=
                \sum_{k=0}^{N-1}a_{k} \,
                        \left( \sum_{n=0}^{N-1} e^{j \, \left( k-r\right) \, \frac{2 \, \pi }{N}\, n}
                        \right) .
\end{equation}

\item Puede probarse que 
\begin{equation}
\sum_{n=0}^{N-1}e^{j \, k \, \frac{2 \, \pi }{N}\, n}=
        \left\{ 
        \begin{array}{l}
                        N, \\ 
                        0,
        \end{array}
        \begin{array}{l}
                k=0,\pm N,\pm 2N,... \\ 
                k\neq 0,\pm N,\pm 2N,...
        \end{array}
        \right. ,
\end{equation}
mediante el siguiente razonamiento
\begin{enumerate}
        \item  Verificamos la igualdad 
               $e^{j \, 0}=e^{j \, N \, \frac{2 \, \pi }{N}\, n}=e^{j \, 2 \, N \, \frac{2 \, \pi}{N}\, n}=...=1.$

        \item  Recordamos la propiedad de las series geom\'{e}tricas convergentes
        \begin{equation}
                \sum_{n=0}^{N-1}a^n=\frac{1-a^{N}}{1-a}, a < 1,
        \end{equation}
        para obtener 
        \begin{equation}
                \sum_{n=0}^{N-1}\left( e^{j \, k \, \frac{2 \, \pi }{N} }\right) ^n=
                        \frac{1-\left( e^{j \, k \, \frac{2 \, \pi }{N} }\right) ^{N}}
                                {1-\left(e^{j \, k \, \frac{2 \, \pi }{N} }\right) }=
                \frac{1-e^{j \, k \, 2 \, \pi }}{1-\left(e^{j \, k \, \frac{2 \, \pi }{N} }\right) }=0.
        \end{equation}
\end{enumerate}

\item Entonces obtenemos 
\begin{equation}
        \sum_{n=0}^{N-1}x\left[ n\right] \, e^{-j \, r \, \frac{2 \, \pi }{N} \, n} = a_{r} \, N,
\end{equation}
o sea 
\begin{equation}
   a_{r} = \frac{1}{N}\sum_{n=0}^{N-1}x\left[ n\right] \, e^{-j \, r \, \frac{2 \, \pi }{N} \, n}. \label{dtfs.ar}
\end{equation}
\end{enumerate}

\subsection{Ejemplo de c\'{a}lculo de los coeficientes de una serie de Fourier}
\begin{enumerate}
\item Consideremos una se\~{n}al peri\'{o}dica con $N = 2$. Supongamos que 
\begin{equation}
x[0] = 1,
\end{equation} y 
\begin{equation}
x[1] = 0.
\end{equation}
Por lo tanto $x[2] = 1$, $x[3] = 0$, y as\'{\i} sucesivamente. \textquestiondown Cuales son los coeficientes correspondientes?

\item Dado que $N = 2$ los coeficientes a calcular son $a_{0}$ y $a_{1}$. Si recordamos la ecuaci\'{o}n (\ref{dtfs.ar}) 
\begin{equation}
   a_{r} = \frac{1}{N}\sum_{n=0}^{N-1}x\left[ n\right] \, e^{-j \, r \, \frac{2 \, \pi }{N} \, n},
\end{equation}
podemos escribir, para $N = 2$,
\begin{equation}
a_{r} = \frac{1}{2} \, 
\left(x[0] \, e^{j \, r \, \frac{2 \, pi}{2} \, 0} + x[1] \, e^{j \, r \, \frac{2 \, pi}{2} \, 1}\right). 
\end{equation}

\item En nuestro caso, para $r = 0$ resulta
\begin{equation}
a_{0} = \frac{1}{2} \, \left( 
x[0] \, e^{j \, 0 \, \frac{2 \, \pi}{2} \, 0} +
x[1] \, e^{j \, 0 \, \frac{2 \, \pi}{2} \, 1} 
\right),
\end{equation}
y para $r = 1$ resulta
\begin{equation}
a_{1} = \frac{1}{2} \, \left( 
x[0] \, e^{j \, 1 \, \frac{2 \, \pi}{2} \, 0} +
x[1] \, e^{j \, 1 \, \frac{2 \, \pi}{2} \, 1} 
\right).
\end{equation}

\item Sustituyendo los valores de $x[0] = 1$ y $x[1] = 0$, obtenemos
\begin{equation}
a_{0} = \frac{1}{2} \, \left( 
1 \times 1 +
0 \times 1
\right) = \frac{1}{2},
\end{equation}
y
\begin{equation}
a_{1} = \frac{1}{2} \, \left( 
1 \times 1 +
0 \times e^{j \, \pi} 
\right) = \frac{1}{2}.
\end{equation}
\end{enumerate}

\subsubsection{Ejercicio}

\begin{enumerate}
\item Calcular los coeficientes de Fourier de una se\~{n}al peri\'{o}dica $N = 2$, donde $x[0] = 0$, y $x[1] = 1$.

\item Calcular los coeficientes de Fourier de una se\~{n}al peri\'{o}dica $N = 3$, donde $x[0] = 0$, $x[1] = 1$ y $x[2] = -1$.
\end{enumerate}

\section{Propiedades de las series de Fourier}

As\'{\i} indicaremos la relaci\'{o}n entre una funci\'{o}n y su serie de
Fourier 
\begin{equation}
x\left[ n\right] 
\begin{array}{l}
FS \\ 
\longleftrightarrow
\end{array}
a_{k}.
\end{equation}

\subsection{Linealidad}

Sean dos funciones con igual per\'{\i}odo $N$
\begin{eqnarray}
x\left[ n\right] &=&x\left[ n+N\right] , \\
y\left[ n\right] &=&y\left[ n+N\right] ,
\end{eqnarray}
con coeficientes de Fourier
\begin{eqnarray}
&&x\left[ n\right] 
\begin{array}{l}
FS \\ 
\longleftrightarrow
\end{array}
a_{k}, \\
&&y\left[ n\right] 
\begin{array}{l}
FS \\ 
\longleftrightarrow
\end{array}
b_{k}.
\end{eqnarray}
Es posible probar
\begin{equation}
z\left[ n\right] = A \, x\left[ n\right] +B \, y\left[ n\right] 
\begin{array}{l}
FS \\ 
\longleftrightarrow
\end{array}
c_{k}=A \, a_{k}+B \, b_{k}.
\end{equation}
Queda como ejercicio para el lector.

\subsection{Desplazamiento en el tiempo}

\begin{enumerate}

\item Dadas las funciones peri\'{o}dicas
\begin{eqnarray}
x\left[ n\right] &=&x\left[ n+N\right] , \\
y\left[ n\right] &=&x\left[ n-n_{0}\right] ,
\end{eqnarray}
sus series de Fourier cumplen con la relaci\'{o}n
\begin{eqnarray}
&&x\left[ n\right] 
\begin{array}{l}
FS \\ 
\longleftrightarrow
\end{array}
a_{k}, \\
&&y\left[ n\right] 
\begin{array}{l}
FS \\ 
\longleftrightarrow
\end{array}
b_{k}=e^{-j \, k \, \frac{2 \, \pi }{N} \, n_{0}} \, a_{k}.
\end{eqnarray}

\item En efecto, escribiendo la ecuaci\'{o}n de an\'{a}lisis
\begin{equation}
b_{k}=
\frac{1}{N} \, \sum_{n=0}^{N-1}x\left[ n-n_{0}\right] \, e^{-j \, k \, \frac{2 \, \pi }{N} \, n},
\end{equation}
donde podemos reemplazar $x[n]$ por su serie de Fourier
\begin{equation}
b_{k}=
\frac{1}{N} \, \sum_{n=0}^{N-1}
\left( \sum_{r=0}^{N-1}a_{r} \, e^{j \, r \, \frac{2 \, \pi}{N} \, \left( n-n_{0}\right) }
\right) \, e^{-j \, k \, \frac{2 \, \pi }{N} \, n}.
\end{equation}

\item Agrupando t\'{e}rminos obtenemos
\begin{equation}
b_{k}=
\frac{1}{N} \, \sum_{n=0}^{N-1}
\left( 
  \sum_{r=0}^{N-1}
  a_{r} \, e^{-j \, r \, \frac{2 \, \pi }{N} \, n_{0}}
  \, e^{j \, \left( r-k\right) \, \frac{2 \, \pi}{N} \, n}
\right).
\end{equation}

\item Es posible agrupar t\'{e}rminos y aplicar la propiedad de suma
  de una serie geom\'{e}trica para obtener
\begin{equation}
\left( \sum_{r=0}^{N-1}
a_{r} \, e^{-j \, r \, \frac{2 \, \pi }{N} \, n_{0}} \, e^{j \, \left( r-k\right) \, \frac{2 \, \pi }{N} \, n}
\right) =
a_{r} \, e^{-j \, r \, \frac{2 \, \pi }{N} \, n_{0}} \, 
\frac{1-e^{j \, \left( r-k\right) \, \frac{2 \, \pi }{N} \, N}}{1-e^{j \, \left( r-k\right) \, \frac{2 \, \pi }{N}}},
\end{equation}

\item Cuando $r\neq k,$ entonces
\begin{equation}
\frac{1-e^{j\left( r-k\right) \frac{2 \, \pi }{N} \, N}}{1-e^{j \, \left(r-k\right) \,
\frac{2 \, \pi }{N}}}
=
\frac{1-e^{j \, \left( r-k\right) \, 2 \, \pi }}{1-e^{j \, \left(
r-k\right) \, \frac{2 \, \pi }{N}}},
\end{equation}
\begin{equation}
\frac{1-e^{j \, \left( r-k\right) \, 2 \, \pi }}{1-e^{j \, \left( r-k\right) \, \frac{2 \, \pi }{N}}}
=
\frac{1-1}{1-e^{j \, \left( r-k\right) \, \frac{2 \, \pi }{N}}}
=0.
\end{equation}

\item Cuando $r=k,$ entonces 
\begin{equation}
\frac{1}{N} \, \sum_{n=0}^{N-1}\left(
a_{k} \, e^{-j \, k \, \frac{2 \, \pi }{N} \, n_{0}} \, 
e^{j \, \left(k-k\right) \, \frac{2 \, \pi }{N} \, n}\right) =
\frac{a_{r} \, e^{-j \, r \, \frac{2 \, \pi }{N} \, n_{0}} \, N}{N},
\end{equation}
y por lo tanto 
\begin{equation}
b_{k}=a_{k} \, e^{-j \, k \, \frac{2 \, \pi }{N} \, n_{0}}.
\end{equation}

\end{enumerate}

\subsection{Desplazamiento en frecuencia}

Consideremos una se\~{n}al per\'{\i}odica $N$, $x\left[ n\right]
=x\left[ n+N\right]$, con coeficientes de Fourier $a_{k}$. 
Es posible probar que al multiplicar $x[n]$ por una arm\'{o}nica $M$, los
coeficientes se desplazan en frecuencia, lo que se expresa como
\begin{equation}
e^{j \, M \, \left( \frac{2 \, \pi }{N}\right) \, n} \, x\left[ n\right] 
\begin{array}{l}
FS \\ 
\longleftrightarrow
\end{array}
a_{k-M}.
\end{equation}

\subsection{Conjugaci\'{o}n}

Consideremos una se\~{n}al per\'{\i}odica $N$, $x\left[ n\right]
=x\left[ n+N\right]$, con coeficientes de Fourier $a_{k}$. 
Es posible probar que al conjugar $x[n]$, los
coeficientes se conjugan y invierten en frecuencia, lo que se expresa como
\begin{equation}
x^{*}\left[ n\right] 
\begin{array}{l}
FS \\ 
\longleftrightarrow
\end{array}
a_{-k}^{*},
\end{equation}

\subsection{Inversi\'{o}n temporal}

Consideremos una se\~{n}al per\'{\i}odica $N$, $x\left[ n\right]
=x\left[ n+N\right]$, con coeficientes de Fourier $a_{k}$. Sea
\begin{equation}
y\left[ n\right] =x\left[ -n\right] ,
\end{equation}

Puede probarse que los coeficientes $a_{k}$ y $b_{k}$ se relacionan
\begin{eqnarray}
y\left[ n\right] 
\begin{array}{l}
FS \\ 
\longleftrightarrow
\end{array}
b_{k}=a_{-k},
\end{eqnarray}
como se puede comprobar
\begin{equation}
x\left[ -n\right] =\sum_{k=0}^{N-1}a_{k} \, e^{-j \, k \, \frac{2 \, \pi }{N} \, n}
=
\sum_{l=1-N}^{0}a_{-l} \, e^{j \, l \, \frac{2 \, \pi }{N}\, n}.
\end{equation}

\subsection{Escalado temporal}

Consideremos una se\~{n}al per\'{\i}odica $N$, $x\left[ n\right]
=x\left[ n+N\right]$, con coeficientes de Fourier
$a_{k}$. Consideremos la se\~{n}al $x_{m} = x[n + m \, N]$,
\begin{equation}
x_{m}\left[ n\right] =\left\{ 
\begin{array}{l}
x\left[ m \right] , \\ 
0
\end{array}
\begin{array}{l}
n=p \, m, \\ 
n=p \, m+q, \, 0 < q < m.
\end{array}
\right.
\end{equation}
Puede probarse la relaci\'{o}n
\begin{equation}
x_{m}\left[ n\right] 
\begin{array}{l}
FS \\ 
\longleftrightarrow
\end{array}
\frac{a_{k}}{m}.
\end{equation}

\subsection{Convoluci\'{o}n peri\'{o}dica}

Consideremos dos se\~{n}ales per\'{\i}dicas $x[n] = x[n + N]$ e $y[n]
= y[n + N]$, con coeficientes de Fourier $a_{k}$ y $b_{k}$. Puede
definirse la operaci\'{o}n convoluci\'{o}n circular y calcularse los
coeficientes de Fourier correspondientes
\begin{equation}
\sum_{r=0}^{N-1} x\left[ r\right] \, y\left[ n-r\right] 
\begin{array}{l}
FS \\ 
\longleftrightarrow
\end{array}
c_{k}=N \, a_{k} \, b_{k}.
\end{equation}

\subsection{Multiplicaci\'{o}n}

Consideremos dos se\~{n}ales per\'{\i}odicas $x[n] = x[n + N]$ e $y[n]
= y[n + N]$, con coeficientes de Fourier $a_{k}$ y $b_{k}$. Puede
definirse la operaci\'{o}n producto y calcularse los
coeficientes de Fourier correspondientes
\begin{equation}
x\left[ n\right] \, y\left[ n\right] 
\begin{array}{l}
FS \\ 
\longleftrightarrow
\end{array}
c_{k}=\sum_{r=0}^{N-1}a_{r} \, b_{k-r}.
\end{equation}

\subsection{Primera diferencia}

Consideremos una se\~{n}al per\'{\i}odica $N$, $x\left[ n\right]
=x\left[ n+N\right]$, con coeficientes de Fourier
$a_{k}$. Puede definirse la operaci\'{o}n primera diferencia y
calcularse los coeficientes de Fourier correspondientes
\begin{equation}
x\left[ n\right] -x\left[ n-1\right] 
\begin{array}{l}
FS \\ 
\longleftrightarrow
\end{array}
\left( 1-e^{-j \, k \, \frac{2 \, \pi }{N}}\right) \, a_{k},
\end{equation}

\subsection{Suma corrida}

Consideremos una se\~{n}al per\'{\i}odica $N$, $x\left[ n\right]
=x\left[ n+N\right]$, con coeficientes de Fourier
$a_{k}$. 
Puede definirse la operaci\'{o}n suma corrida y
calcularse los coeficientes de Fourier correspondientes
\begin{equation}
y\left[ n\right] =\sum_{k=-\infty }^nx\left[ k\right] 
\begin{array}{l}
FS \\ 
\longleftrightarrow
\end{array}
 \frac{a_{k}}{\left( 1-e^{-j \, k \, \frac{2 \, \pi }{N}}\right) }.
\end{equation}

\subsection{Conjugaci\'{o}n y simetr\'{\i}a de la conjugaci\'{o}n}

Consideremos una se\~{n}al per\'{\i}odica $N$, $x\left[ n\right]
=x\left[ n+N\right]$, con coeficientes de Fourier
$a_{k}$. Entonces pueden verificarse la simetr\'{\i}a de la conjugaci\'{o}n
\begin{equation}
a_{-k}=\overline{a_{k}}.
\end{equation}

\subsection{Se\~{n}ales reales pares}

Consideremos una se\~{n}al per\'{\i}odica $N$, $x\left[ n\right]
=x\left[ n+N\right]$, con coeficientes de Fourier
$a_{k}$. La se\~{n}al $x[n]$ es par si y solo si la parte imaginaria
de sus coeficientes es nula para todo k, o en otras palabras,
\begin{equation}
a_{-k}=a_{k}.
\end{equation}

\subsection{Se\~{n}ales reales impares}

Consideremos una se\~{n}al per\'{\i}odica $N$, $x\left[ n\right]
=x\left[ n+N\right]$, con coeficientes de Fourier
$a_{k}$. La se\~{n}al $x[n]$ es impar si y solo si la parte real
de sus coeficientes es nula para todo k, o en otras palabras,
\begin{equation}
a_{-k} = -a_{k}.
\end{equation}

\subsection{Parte par e impar de se\~{n}ales reales}

La partes pares e impares de una se\~{n}al real de tiempo discreto
$x[n] = x[n + N]$ con coeficientes de Fourier $a_{k}$ pueden calcularse a partir de sus coeficientes de Fourier
\begin{eqnarray}
x\left[ n\right] &=&x_{p}\left[ n\right] +x_i\left[ n\right] , \\
x_{p}\left[ n\right] &=&x_{p}\left[ -n\right] , \\
x_i\left[ n\right] &=&-x_i\left[ -n\right] , \\
x_{p}\left[ n\right] 
&\begin{array}{l}
FS \\ 
\longleftrightarrow
\end{array}&
\Re\left( a_{k}\right) , \\
x_i\left[ n\right] 
&\begin{array}{l}
FS \\ 
\longleftrightarrow
\end{array}&
\Im\left( a_{k}\right).
\end{eqnarray}

\subsection{Relaci\'{o}n de Parseval}

\begin{enumerate}

\item La potencia media de una se\~{n}al peri\'{o}dica $N$ es
\begin{equation}
        \frac{1}{N}\sum_{n=0}^{N-1}\left| x\left[ n\right]
        \right|^{2}=
        \frac{1}{N}\sum_{n=0}^{N-1} x\left[ n\right] \, \overline{x\left[ n \right]}
\end{equation}

\item Sustituyendo la funci\'{o}n conjugada por su serie de Fourier obtenemos
\begin{equation}
        \frac{1}{N}\sum_{n=0}^{N-1}\left| x\left[ n\right]
        \right|^{2}=
        \frac{1}{N} \sum_{n=0}^{N-1} x\left[ n\right] \, 
        \left(\frac{1}{N}\sum_{k=0}^{N-1} \overline{a_{k}} \, e^{-j \, \frac{2 \, \pi}{N} \, k \, n} \right).
\end{equation}

\item Intercambiando sumatorias obtenemos
\begin{equation}
        \frac{1}{N}\sum_{n=0}^{N-1}\left| x\left[ n\right]
        \right|^{2}=
        \frac{1}{N}\sum_{k=0}^{N-1} \overline{a_{k}}  \, \left(
        \frac{1}{N}\sum_{k=0}^{N-1} x\left[ n\right]  \, e^{-j \, \frac{2 \, \pi}{N} \, k \, n} \right).
\end{equation}

\item Finalmente, observemos que la expresi\'{o}n entre par\'{e}ntesis
  puede sustituirse por $a_{k}$
\begin{equation}
        \frac{1}{N}\sum_{n=0}^{N-1}\left| x\left[ n\right]
        \right|^{2}=
        \sum_{k=0}^{N} a_{k} \, \overline{a_{k}}.
\end{equation}

\item La secuencia $|a_{k}|^{2}$, $k = 0, 1, ... N - 1$ se denomina Espectro
de Densidad de Potencia.

\end{enumerate}

\section{C\'{a}lculo de los coeficientes, otra vez}

\subsection{FFT, Teoría}

A continuación vamos a ver el desarrollo de ideas de la transformada rápida de Fourier, o 
Fast Fourier Transform.

\begin{enumerate}
\item Si observamos la f\'{o}rmula para calcular los coeficientes de las series de Fourier de tiempo discreto,
\begin{equation}
   a_{r} =  \frac{1}{N} \, \sum_{n=0}^{N-1}x\left[ n\right] \, e^{-j \, r \, \frac{2 \, \pi}{N} \, n},
\end{equation}
notaremos que es posible expresarla como
\begin{equation}
   a_{r} =  \frac{1}{N} \, \sum_{n=0}^{N-1}x\left[ n\right] \, W^{n \, r},
\end{equation}
definiendo
\begin{equation}
  W = e^{-j \, \frac{2 \, \pi}{N}}.
\end{equation}

\item Con esta transformaci\'{o}n es inmediato ver que el coeficiente $a_{r}$ puede expresarse como el producto de dos vectores
\begin{equation}
N \, a_{r} =
\left[\begin{array}{ccccc}
W^{r \, 0} & ... & W^{r \, n} & ... & W^{r \, (N-1)}
\end{array} \right] \cdot 
\left[\begin{array}{c}
x[0]\\
... \\
x[n]\\
... \\
x[N-1]
\end{array} \right].
\end{equation}

\item Por lo tanto, todos los coeficientes de la serie de Fourier pueden calcularse como un producto de 
matriz por vector
\begin{equation}
N \, \left[ \begin{array}{c}
a[0] \\
... \\
a[r] \\
... \\
a[N-1]
\end{array} \right] =
\left[ \begin{array}{cccccc}
W^{0 \, 0} & ... & W^{0 \, n} & ... & W^{0 \, (N-1)} \\
... & .... & ... & ... & ... \\
W^{r \, 0} & ... & W^{r \, n} & ... & W^{r \, (N-1)} \\
... & .... & ... & ... & ... \\
W^{(N-1) \, 0} & ... & W^{(N-1) \, n} & ... & W^{(N-1) \, (N-1)}
\end{array} \right] =
\left[\begin{array}{c}
x[0]\\
... \\
x[n]\\
... \\
x[N-1]
\end{array} \right].
\end{equation}

\item Esta \'{u}ltima ecuaci\'{o}n transforma 
la representaci\'{o}n $x[n]$ del dominio del tiempo en 
la representaci\'{o}n $a[k]$ del dominio de la frecuencia.

\item Podemos definir entonces la matriz
\begin{equation}
\mathbf{M}_{N} = 
\left[ \begin{array}{cccccc}
W^{0 \, 0} & ... & W^{0 \, n} & ... & W^{0 \, (N-1)} \\
... & .... & ... & ... & ... \\
W^{r \, 0} & ... & W^{r \, n} & ... & W^{r \, (N-1)} \\
... & .... & ... & ... & ... \\
W^{(N-1) \, 0} & ... & W^{(N-1) \, n} & ... & W^{(N-1) \, (N-1)}
\end{array} \right].
\end{equation}

\item La matriz $\mathbf{M}_{N}$ tiene propiedades interesantes:
\begin{enumerate}
\item La matriz $\mathbf{M}_{N}$ es simétrica.
\item El producto de $\mathbf{M}_{N}$ por su transpuesta es igual a $N$ veces la matriz identidad 
$\mathbf{I}$
\begin{equation}
  \mathbf{M}_{N}^{t} \cdot \mathbf{M}_{N} = \mathbf{M}_{N} \cdot \mathbf{M}_{N}^{t} = N \, \mathbf{I}.
\end{equation}
\item El vector de coeficientes de Fourier $\mathbf{a}$
puede obtenerse con la fórmula
\begin{equation}
 N \, \mathbf{a} = \mathbf{M}_{N} \cdot \mathbf{x},
\end{equation}
y el vector $x$ muestras en el dominio del tiempo puede recuperarse a partir de los coeficientes 
$\mathbf{a}$
\begin{equation}
 N \, \mathbf{M}_{N}^{t} \cdot \mathbf{a} = \mathbf{x}.
\end{equation}


\end{enumerate}


\end{enumerate}


\subsection{FFT, práctica}

\subsubsection{Como generar la matriz M}

Generamos la matriz $M$ y verificamos sus propiedades.

\lstinputlisting{./fft/demo_matriz.m}

\subsection{Como usar la función fft de Matlab/Octave}

\lstinputlisting{./fft/myfft.m}

\lstinputlisting{./fft/demofft.m}

\lstinputlisting{./fft/demofase.m}

\subsection{Las diferencias entre notas musicales de diferentes instrumentos}

\lstinputlisting{./fft/demo_instrumentos.m}
